\uuid{UOKy}
\exo7id{3706}
\titre{exo7 3706}
\auteur{quercia}
\organisation{exo7}
\datecreate{2010-03-11}
\isIndication{false}
\isCorrection{true}
\chapitre{Espace euclidien, espace normé}
\sousChapitre{Espace vectoriel euclidien de dimension 3}
\cpge{mpsi-mp2i.applications-lineaires.a, pcsi.applications-lineaires.a, ptsi.applications-lineaires.a, mpsi-mp2i.espaces-prehilbertiens.d, pcsi.espaces-prehilbertiens.d, bcpst1.applications-lineaires-matrices, mp-mpi.endomorphismes-euclidien.b, pc.endomorphismes-euclidiens.b, psi.endomorphismes-euclidien.b, pt.isometries-euclidien.b, mp-mpi.endomorphismes-euclidien.c, pc.endomorphismes-euclidiens.a, psi.endomorphismes-euclidien.a, pt.isometries-euclidien.a, bcpst2.applications-lineaires-matrices.a}
\module{Algèbre}
\niveau{L2}
\difficulte{}

\contenu{
\texte{
Soit $E$ un espace vectoriel euclidien et $f \in \mathcal{L}(E)$ antisymétrique.
}
\begin{enumerate}
    \item \question{Montrer que $\mathrm{id}_E + f \in GL(E)$.}
    \item \question{Montrer que $g = (\mathrm{id} - f) \circ (\mathrm{id} + f)^{-1} \in {\cal O}^+(E)$ et
    $\mathrm{id}+g$ est inversible.}
    \item \question{Réciproquement, soit $h \in {\cal O}^+(E)$ tq $\mathrm{id} + h$ soit inversible.
    Montrer qu'il existe $f$ antisymétrique tel que
    $h = (\mathrm{id} - f) \circ (\mathrm{id} + f)^{-1}$.}
\reponse{
matrice dans une {\it bond}.
$f = 2(h+\mathrm{id})^{-1} - \mathrm{id}$.
}
\end{enumerate}
}
